%%%%%%%%%%%%%%%%%%%%%%%%%%%%%%%%%%%%%%%%%%%%%%%%%%%%%%%%%%%%%%%%%
\chapter{LİTERATÜR ÖZETİ}\label{ch:literatur}
%%%%%%%%%%%%%%%%%%%%%%%%%%%%%%%%%%%%%%%%%%%%%%%%%%%%%%%%%%%%%%%%%
\vspace*{-1cm}

Kişisel finans yönetimi (PFM - Personal Finance Management), bireylerin gelir ve giderlerini takip ederek mali durumlarını optimize etmelerini sağlayan kritik bir süreçtir. Bu bölümde, finansal teknolojilerin (FinTech) gelişimi, mevcut uygulamalar ve yapay zekânın bu alandaki rolü literatür ışığında incelenmiştir.

\section{Kişisel Finans Yönetiminin Gelişimi}

Finansal takip süreçleri, geleneksel olarak kağıt üzerindeki kayıtlardan dijital tablolara (Excel vb.) ve sonrasında modern mobil/web uygulamalarına evrilmiştir. Mint, YNAB (You Need A Budget) ve PocketGuard gibi popüler platformlar, kullanıcılara harcamalarını kategorize etme ve bütçe hedefleri belirleme imkanı sunmaktadır. Ancak, bu sistemlerin çoğu hala yoğun manuel veri girişi gerektirmekte ve geleceğe yönelik dinamik rehberlik sunmakta yetersiz kalmaktadır \cite{moore91}.

\section{Yapay Zekâ ve Finansal Öngörü}

Yapay zekâ (YZ), finans dünyasında dolandırıcılık tespiti, kredi skorlama ve portföy yönetimi gibi alanlarda uzun süredir kullanılmaktadır. Kişisel finans tarafında ise, harcama alışkanlıklarının analizi için makine öğrenmesi modelleri giderek daha fazla önem kazanmaktadır.

\subsection{Zaman Serisi Analizi ve LSTM}
Finansal veriler, doğası gereği zaman serisi niteliğindedir. Geleneksel istatistiksel yöntemler (ARIMA vb.) karmaşık ve doğrusal olmayan harcama kalıplarını yakalamakta zorlanabilir. Uzun Kısa Süreli Bellek (LSTM) ağları, geçmiş verilerdeki uzun vadeli bağımlılıkları öğrenebildiği için harcama tahminleme süreçlerinde yüksek başarı göstermektedir \cite{Wegener2000629}.

\subsection{Dürtme (Nudge) Teorisi ve Bütçeleme}
Davranışsal iktisatın bir parçası olan "Dürtme Teorisi", bireylerin özgür iradelerini kısıtlamadan onları daha iyi kararlar almaya yönlendirmeyi amaçlar. Yapay zekâ tabanlı bütçe önerileri, kullanıcıya "ideal" harcama hedefleri sunarak bir dijital dürtme mekanizması işlevi görür. Hibrit modeller, kullanıcının gerçek hayat koşulları ile ideal finansal hedefleri arasında denge kurarak bu dürtmelerin etkisini artırır.

\section{OCR ve Veri Giriş Otomasyonu}

Kullanıcıların harcama takip uygulamalarından vazgeçmelerinin en büyük nedeni olan manuel veri girişi yükü, Optik Karakter Tanıma (OCR) teknolojileri ile aşılmaya çalışılmaktadır. Tesseract gibi açık kaynaklı OCR motorları, mobil cihazlarla çekilen makbuz fotoğraflarından tarih, miktar ve satıcı bilgilerini ayıklayarak süreci otomatize etmektedir \cite{harper2007}.

Sonuç olarak literatür, kişisel finans yönetiminde veri girişi kolaylığı ve yapay zekâ destekli proaktif yönlendirmelerin, bireylerin finansal sağlığını iyileştirmede en kritik faktörler olduğunu göstermektedir.
